\documentclass[letterpaper,10pt]{article}
\usepackage[left=2cm,top=3cm,right=2cm,bottom=4cm]{geometry} 
\usepackage[utf8]{inputenc}
\usepackage[activeacute,spanish]{babel}
\usepackage{multirow}	% Para poder unir filas en las tablas
\usepackage{colortbl}	% Para colorear tablas
\usepackage{latexsym,amssymb,oldlfont,amsmath}
\usepackage{graphicx,graphics}
\usepackage{multicol}
\usepackage{color}
\usepackage{listings}
\usepackage{lastpage}
\usepackage{enumitem}
\usepackage{mathtools}



%\usepackage{fontspec}
%\setmainfont{Times New Roman}
\newcommand{\R}{{\mathbb R}}
\newcommand{\si}{{\mathcal S}^n}
\renewcommand{\theequation}{\thesection.\arabic{equation}} % Cambia formato salida numeracion equation
\makeatletter\@addtoreset {equation}{section} \makeatother % Tambien se usa
\renewcommand{\baselinestretch}{1.3}
\usepackage{amssymb}
\usepackage{hyperref}
\definecolor{midnightblue}{rgb}{0.1,0.1,0.44}


\newcommand{\dis}{\displaystyle}
\renewcommand{\theenumii}{\arabic{enumii}}

\def\Z{\mathbb{Z}}
\def\N{\mathbb{N}}
\def\R{\mathbb{R}}
\def\M{\mathcal{M}_n(\mathbb{R})}
\def\sin{\operatorname{sen}}
\def\e{\operatorname{e}}
\DeclareMathOperator{\arcsec}{arcsec}
\DeclareMathOperator{\arccot}{arccot}
\DeclareMathOperator{\arccsc}{arccsc}
\providecommand{\abs}[1]{\lvert#1\rvert}


\usepackage{etoolbox,fancyhdr,xcolor}
\fancyhf{} % sets both header and footer to nothing
\renewcommand{\headrulewidth}{0pt}
\pagestyle{fancy}
\headheight=60pt %para cambiar el tamaño del encabezadox
\fancyhead[L] %la "L" indica a la izquierda
{
\begin{minipage}{6cm}
\includegraphics[width=0.9\textwidth]{logo_udp.eps}
\end{minipage}
}


\fancyfoot[C]{}

\title{C2_Ay6_pauta}

\begin{document}

\begin{center}
Cálculo II\\
Ayudantía Nº6 - Pauta\\
Primer Semestre 2017
\end{center}

\begin{enumerate}
\item Determine si las siguientes afirmaciones son verdaderas o falsas y justifique su respuesta:

\begin{enumerate}[label=\emph{\alph*)}]

\item Si $h$ es una función continua y $f$, $g$ son derivables, entonces para la función:\\
\begin{center}
$\dis F(x) = \int_{f(x)}^{g(x)} h(t)\, dt$
\end{center} 
Se tiene que $F'(x)=h(g(x))g'(x)-h(f(x))g'(x)$

Aplicando el teorema fundamental del cálculo a $F(x)$ se obtiene lo siguiente:

\begin{center}
$\dis F(x) = \int_{f(x)}^{g(x)} h(t)\, dt$\\
$F'(x)=h(g(x))g'(x)-h(f(x))f'(x)$
\end{center} 

Por lo tanto la afirmación es \underline{FALSA}.

\dotfill\textbf{[1 punto]}


\item Si $f(x)$ es una función continua y positiva entonces:
\begin{center}
$\dis F(x) = \frac{\int_{0}^{x} t * f(t)\, dt}{\int_{0}^{x} f(t)\, dt}$
\end{center} 
Es creciente para $x \geq 0$\\

\begin{align*}
 \dis  F(x)  &= \frac{\int_{0}^{x} t * f(t)\, dt}{\int_{0}^{x} f(t)\, dt}\\
 \dis  F'(x)  &= \frac{x*f(x)\int_{0}^{x} f(t)\, dt - f(x)\int_{0}^{x} t * f(t)\, dt}{(\int_{0}^{x} f(t)\, dt)^2}\\
 \dis &= \frac{f(x)(x\int_{0}^{x} f(t)\, dt - \int_{0}^{x} t * f(t)\, dt)}{(\int_{0}^{x} f(t)\, dt)^2}\\
 \dis &= \frac{f(x)(\int_{0}^{x} xf(t)\, dt - \int_{0}^{x} t * f(t)\, dt)}{(\int_{0}^{x} f(t)\, dt)^2}\\
 \dis &= \frac{f(x)(\int_{0}^{x} (x-t)f(t)\, dt)}{(\int_{0}^{x} f(t)\, dt)^2} \geq 0, \forall x \geq 0
\end{align*}

Notar que $\displaystyle \int_{0}^{x} (x-t)f(t)\, dt$ es positivo ya que $0 \leq t \leq x$ y $f(t)>0$.

Por lo tanto la afirmación es \underline{VERDADERA}.

\dotfill\textbf{[1 punto]}



\item Si $f$ es continua y $\int_{0}^{4} f(x)\, dx = 40$, entonces $\int_{0}^{2} f(2x)\, dx = 5$\\

Utilizando la sustitución $u = 2x$ y $du = 2 dx $

\begin{align*}
 \dis  \int_{0}^{2} f(2x)\, dx  &= \frac{\int_{0}^{4} f(u)}{2}\, dx\\
 \dis    &= \frac{1}{2}\int_{0}^{4} f(u)\, du\\
 \dis    &= \frac{1}{2}* 40 = 20\\
\end{align*}

Por lo tanto la afirmación es \underline{FALSA}.

\dotfill\textbf{[1 punto]}

\end{enumerate}
\vspace{2mm}







\item Calcule:
$\displaystyle \lim_{x\to 0} \; {1\over x} \; \int ^x_0 (1 + \sin
2t)^{{1\over t}} dt$

\begin{align*}
 \dis  \displaystyle \lim_{x\to 0} \; {1\over x} \; \int ^x_0 (1 + \sin 2t)^{{1\over t}} dt  &= \displaystyle \lim_{x\to 0} \; \frac{ \int ^x_0 (1 + \sin 2t)^{{1\over t}} dt}{x} = \; {0\over 0}\\
 \dis    &= \displaystyle \lim_{x\to 0} \frac{(1 + \sin 2x)^{{1\over x}}}{1}\\
 \dis    &= \displaystyle \lim_{x\to 0}  (1 + \sin 2x)^{{1\over x}}\\
 \dis    &= \displaystyle \lim_{x\to 0}  (1 + \frac{\sin 2x}{2x} 2x)^{{1\over x}}\\
 \dis    &= e^2
\end{align*}

El valor obtenido proviene del hecho que $\displaystyle \lim_{x\to 0} \frac{\sin(2x)}{2x}=1$ y $\displaystyle \lim_{x\to 0} \left(1+2x\right)^{\frac{1}{x}} = e^2$

\dotfill\textbf{[1 punto]}

\newpage








\item Si $f$ es continua y $\displaystyle \int_0^9f(x)dx=4$, encuentre $\displaystyle \int_0^3xf(x^2)dx$\\

Utilizando la sustitución $u = x^2$ y $du = 2x \, dx$ obtenemos lo siguiente:

\begin{align*}
 \dis  \displaystyle \int_0^3xf(x^2)dx  &= \displaystyle \int_0^9f(u) \frac{du}{2}\\
 \dis    &= \frac{1}{2}\int_{0}^{9} f(u)\, du\\
 \dis    &= \frac{1}{2}* 4 = 2\\
\end{align*}

\dotfill\textbf{[1 punto]}




\item Calcule la siguiente integral definida: $\dis \int_{0}^{\pi/6} \frac{\sin(2x)}{\cos^2(2x)} \; dx$

Utilizando la sustitución $u = \cos (2x)$ y $du = -2\sin (2x) \, dx$ obtenemos lo siguiente:

\begin{align*}
 \dis \int_{0}^{\pi/6} \frac{\sin(2x)}{\cos^2(2x)} \; dx  &= \displaystyle -\frac{1}{2} \int_1^{\frac{1}{2}} \frac{du}{u^2}\\
 \dis    &= \displaystyle \frac{1}{2} \int_{\frac{1}{2}}^1 \frac{du}{u^2}\\
 \dis    &= -\frac{1}{2} \left.\frac{1}{u}\right|_{\frac{1}{2}}^1 = \frac{1}{2}
\end{align*}

\dotfill\textbf{[1 punto]}


\item Calcule la siguiente integral definida: $\dis \int_{1/2}^{\sqrt{3}/2} \frac{6}{\sqrt{1-t^2}} \; dt$

Utilizando la sustitución $t = \sin (x)$ y $dt = \cos (x) \, dx$ obtenemos lo siguiente:

\begin{align*}
 \dis \int_{1/2}^{\sqrt{3}/2} \frac{6}{\sqrt{1-t^2}} \; dt  &= \dis \int_{\pi/6}^{\pi/3} \frac{6}{\sqrt{1-\sin^2x}} \; \cos x \; dx\\
 \dis    &= \dis \int_{\pi/6}^{\pi/3} 6 \; dx\\
 \dis    &=\displaystyle  \left.6x\right|_{\pi/6}^{\pi/3} = \pi
\end{align*}

\dotfill\textbf{[1 punto]}

\end{enumerate}




\vspace{10mm}




\end{document}
